\documentclass[10pt,a4paper]{amsart}
\usepackage[margin=19mm]{geometry}
\usepackage{amsmath,amssymb,mathtools}
\usepackage[scr=rsfs,cal=euler]{mathalfa}
\usepackage{xfrac}
\usepackage[unicode,hidelinks,bookmarks=false]{hyperref}
\newcommand{\ie}{i.e.\ }
\newcommand{\cf}{cf.\ }
\newcommand{\resp}{respectively\ }
\newcommand{\Subsection}{\subsection}
\providecommand{\nobreakdash}{\nobreak\hbox{-}}
\setlength{\emergencystretch}{2em}
\multlinegap0pt
\usepackage{amsthm}
\usepackage{enumitem}
\providecommand{\bysame}{\textemdash\textemdash\textemdash\ }
\providecommand{\newblock}{\space}
\newlist{thmlist}{enumerate}{1}
\setlist[thmlist]{label=(\roman{thmlisti}),noitemsep}
\newtheorem{theorem}{Theorem}[section]
\newtheorem{proposition}[theorem]{Proposition}
\newtheorem{lemma}[theorem]{Lemma}
\newtheorem{Claim}[theorem]{Claim}
\newtheorem{claim}[theorem]{Claim}
\newtheorem{fact}[theorem]{Fact}
\newtheorem{problem}[theorem]{Problem}
\newtheorem*{claim*}{Claim}
\newtheorem{corollary}[theorem]{Corollary}
\newtheorem{Main Conjecture}[theorem]{Main Conjecture}
\newtheorem{conjecture}[theorem]{Conjecture}
\theoremstyle{remark}
\newtheorem{defn}[theorem]{Definition}
\newtheorem{examplex}[theorem]{Example}
\newenvironment{example}
  {\pushQED{\qed}\renewcommand{\qedsymbol}{$\diamondsuit$}\examplex}
  {\popQED\endexamplex}
\newtheorem{Remark}[theorem]{Remark}
\newenvironment{remark}
  {\pushQED{\qed}\renewcommand{\qedsymbol}{$\diamondsuit$}\remarkx}
  {\popQED\endremarkx}
\newtheorem{remarkx}[theorem]{Remark}
\newtheorem{notation}[theorem]{Notation}
\theoremstyle{plain}
\newtheorem{question}{Question}
\newtheorem*{MainThm}{Main Theorem}
\newtheorem*{Fact}{Fact}
\newtheorem*{conj}{Conjecture}
\newcommand{\mydef}[1]{{\bf #1}}
\newcommand{\notdefterm}[1]{\emph{#1}}
\newcommand{\inctabs}[2]{{\rm Inc}(#1,#2)}
\newcommand{\groth}[1]{{\mathfrak G}_{#1}}
\newcommand{\symgroth}[1]{{g}_{#1}}
\newcommand{\stabgroth}[1]{G_{#1}}


\newcommand{\init}{{\tt in}}
\newcommand{\mult}{{\tt mult}}
\newcommand{\bpd}{{\sf BPD}}
\newcommand{\pipes}{{\sf Pipes}}
\newcommand{\hgt}{{\rm{ht }}}


\newcommand{\ess}{{\tt Ess}}
\newcommand{\perm}{{\tt Perm}}
\newcommand{\asm}{{\tt ASM}}
\newcommand{\wt}{{\tt wt}}

\newcommand{\ZZ}{\mathbb{Z}}
\newcommand{\CC}{\mathbb{C}}
\newcommand{\NN}{\mathbb{N}}
\newcommand{\QQ}{\mathbb{Q}}

\newcommand{\hilb}{\mbox{Hilb}}

\newcommand{\rk}{{\tt rk}}

\newcommand{\mat}{{\sf Mat}}
\newcommand{\spec}{{\mbox{Spec}}}
\newcommand{\proj}{{\mbox{Proj}}}

\newcommand{\fl}{{\mathcal F}}
\newcommand{\GL}{{\sf GL}}

%notation for taking w to v when doing transition.  Put subscripts with location of lower outside corner.
\newcommand{\vmap}{\xi}
\title[Bumpless pipe dreams encode Gr"obner geometry of Schubert polynomials]{Bumpless pipe dreams encode Gr"obner geometry of Schubert polynomials: the ASM ideal facts (Lemma 2.6)}
\author{Patricia Klein and Anna Weigandt}
\date{Source: arXiv:2108.08370v4 (CC BY 4.0)}
\begin{document}
\maketitle
\noindent\emph{Source and licence.} Bumpless pipe dreams encode Gr"obner geometry of Schubert polynomials, arXiv:2108.08370v4, distributed under the Creative Commons Attribution 4.0 International licence, http://creativecommons.org/licenses/by/4.0/, as recorded in the arXiv OAI metadata for this exact version and shown on the abstract page https://arxiv.org/abs/2108.08370v4 (v4 of 2 September 2025; version v3 of this e-print was withdrawn by the authors and is not used here). Full rights notice: licenses/arxiv-2108.08370-CC-BY-4.0.txt.\par\smallskip
\noindent\textit{Editorial scope.} Counted unit: the six-part lemma on ASM ideals, source label \textbackslash{}label\{lemma:asmIdealFacts\} (source0.tex lines 323--342), printed as Lemma 2.6 of the article, delivered with its complete original proof, the whole of the authors\textquotesingle{} \textbackslash{}begin\{proof\} block at source0.tex lines 344--384 (parts (i)--(vi), from the containment chain eq:idealcontainments to the two-line conclusion of part (vi)), as the single primary proof body.\par
\noindent\textit{Required context retained uncounted.} Two source spans: the statement of Proposition 2.3 on matrix Schubert facts (source0.tex lines 283--291, source label \textbackslash{}label\{prop:matSchubFacts\}), which the proof invokes in part (iv); and the statement of the Knutson--Miller theorem (source0.tex lines 310--315, source label \textbackslash{}label\{theorem:KMTheoremB\}), printed as Theorem 2.4, which the proof invokes in part (i). The three spellings \textbackslash{}ref\{lempart:addIwi\}, \textbackslash{}ref\{lempart:asmDecomp\} and \textbackslash{}ref\{lempart:asmHeight\} inside the proof resolve against the item labels that the lemma statement itself carries in the delivered text.\par
\noindent\textit{One disclosed adaptation of a cross-reference.} The proof\textquotesingle s phrase (see Subsection \textbackslash{}ref\{subsect:hilb\}) points at a subsection of the article that lies wholly outside this unit; the single command \textbackslash{}ref\{subsect:hilb\} is therefore delivered as \textbackslash{}href\{https://arxiv.org/abs/2108.08370v4\}\{2.4\}, a hyperlink to the version of record whose visible text is the subsection number that the version PDF prints at that place. This is the only replacement in the whole package; it is recorded as a reference-only adaptation in the delivery evidence, and the surrounding words, the sentence and every other symbol of that proof line are unchanged.\par
\noindent\textit{Wrapper and numbering.} The article\textquotesingle s \textbackslash{}documentclass[12pt,reqno]\{amsart\} and its package list are replaced by the AMS amsart wrapper of the pinned TeX Live runtime (10pt a4 with the wrapper\textquotesingle s geometry); the authors\textquotesingle{} own theorem environments (source lines 36--63) and their own macros (source lines 67--106) are carried over verbatim, together with the enumitem list environment thmlist that the lemma statement uses. The authors\textquotesingle{} drawing/tableau packages (tikz, xy, ytableau, bbold, wrapfig) are not carried over: no drawing, tableau or graphical construct lies in the delivered spans, and the 47 tikzpictures of the article are not redistributed. Editorial counter commands placed between bodies restore the article\textquotesingle s own printed numbers: Proposition 2.3 and Theorem 2.4 for the two context statements and Lemma 2.6 for the counted statement, each of which is the number the version PDF prints (page 9). The article\textquotesingle s section-based equation numbering is not reproduced; the single numbered display of the proof (eq:idealcontainments) carries the wrapper\textquotesingle s own number, and every internal cross-reference resolves to the correct object inside the unit.\par
\noindent\textit{Citation map.} The delivered unit invokes exactly six citations: \textbackslash{}cite[Proposition 3.3]\{Ful92\} in the context proposition, \textbackslash{}cite[Theorem B]\{KM05\} in the context theorem, and \textbackslash{}cite[Proposition 4.8]\{Wei\}, \textbackslash{}cite{LS96}, \textbackslash{}cite[Proposition 9]\{Rea02\} and \textbackslash{}cite[Lemma 15.19]\{MS04\} in the counted proof. Their six entries are transcribed from the article\textquotesingle s own external bibliography main.bbl, kept verbatim in eprint/main.bbl in this item directory, and print with the article\textquotesingle s own author-year labels [Ful92], [KM05], [Wei17], [LS96], [Rea02] and [MS04]. The LS96 entry begins with the source\textquotesingle s own \textbackslash{}bysame, which in the article refers to the immediately preceding bibliography entry; that neighbour is not cited by this unit and is therefore not delivered, and the dash is kept literal rather than silently replaced by author names.\par
\noindent\textit{Assets.} The e-print of this version ships the author .tex, the external bibliography main.bbl and a 00README.json; it contains no class, style, font or bitmap. No retained span contains any \textbackslash{}includegraphics, \textbackslash{}input or \textbackslash{}include command; the article\textquotesingle s 47 tikzpictures and its one figure environment lie outside the delivered spans and none of them is redistributed. Every mathematical symbol, sentence and label of the delivered spans is the authors\textquotesingle{} own native text; no mathematics was written, repaired, re-derived or normalized, and no printed slip inside these spans was corrected.\par
\setcounter{section}{2}
\setcounter{theorem}{2}
\begin{proposition}[{\cite[Proposition~3.3]{Ful92}}]
\label{prop:matSchubFacts}
	Fix $w\in S_n$.
	\begin{enumerate}
		\item $I_w$ is prime.
		\item $\hgt (I_w)=\ell(w)$.
		\item $X_w$ is Cohen--Macaulay.
	\end{enumerate}
\end{proposition}

\setcounter{theorem}{3}
\begin{theorem}[{\cite[Theorem~B]{KM05}}]
\label{theorem:KMTheoremB}
If $\sigma$ is any anti-diagonal term order on $\kappa[z_{1,1}, \ldots, z_{n,n}]$, then the Fulton generators of $I_w$ 
form a Gr\"obner basis.  In particular,  $\init_\sigma(I_w)$ is squarefree.  The facets of $\Delta(\init_\sigma(I_w))$ are \[\{[n]\times[n] - C(\mathcal D):\mathcal D\in \pipes(w)\},\]
and \[\init_\sigma(I_w)=\bigcap_{\mathcal D\in \pipes(w)} I_{C(\mathcal D)}.\]
\end{theorem}

\setcounter{theorem}{5}
\begin{lemma}\label{lemma:asmIdealFacts}
Let $A\in \asm(n)$, and fix an anti-diagonal term order $\sigma$ on $ R = \kappa[z_{1,1},\ldots,z_{n,n}]$.   Then the following hold:
\begin{thmlist}
\item  If $w_1,\ldots,w_r\in S_n$ such that $A=\vee\{w_1,\ldots,w_r\}$, then
\[ \sum_{i=1}^r \init_\sigma(I_{w_i})=\init_\sigma(I_A)=\bigcap_{u\in \perm(A)} \init_\sigma(I_u).\]
\label{lempart:intersectInit}
    

     
    \item If $w_1,\ldots,w_r\in S_n$ such that $A=\vee\{w_1,\ldots,w_r\}$, then $ I_A={\displaystyle \sum_{i=1}^r} I_{w_i}.$ \label{lempart:addIwi}
    
     \item $I_A$ is radical. \label{lempart:asmRad}
    
   \item $I_A$ has the irredundant prime decomposition $\displaystyle I_A=\bigcap_{w\in \perm(A)}I_w$. \label{lempart:asmDecomp}
   
\item $\hgt (I_A)=\hgt(A)$. \label{lempart:asmHeight}

		\item $A$ is equidimensional if and only if $\spec(R/I_A)$ is equidimensional. \label{lempart:asmEquidim}
\end{thmlist}
\end{lemma}

\begin{proof}

\noindent (i) Since $A\geq w_{i}$  we have $\rk_A\leq \rk_{w_i}$, for all $i\in [r]$.  Thus, $I_{w_i}\subseteq I_{A}$ and so \[\sum_{i=1}^r I_{w_i}\subseteq I_A.\]
Similarly, since $A\leq u$ for all $u\in \perm(A)$, we know 
$ I_A\subseteq \bigcap\{ I_u:u\in \perm(A)\}$.

As such,
\begin{equation}
\label{eq:idealcontainments}
    \sum_{i=1}^r \init_\sigma(I_{w_i})\subseteq \init_\sigma\left(\sum_{i=1}^r I_{w_i}\right)\subseteq \init_\sigma(I_A)\subseteq \init_\sigma\left(\bigcap_{u\in \perm(A)}I_{u}\right)\subseteq \bigcap_{u\in \perm(A)}\init_\sigma(I_u).
\end{equation}

By Theorem~\ref{theorem:KMTheoremB}, for all $w\in S_n$, $\init_\sigma(I_w)$ is a squarefree monomial ideal with associated Stanley-Reisner complex $\Delta(\init_\sigma(I_w))$.  In particular, this implies that 
$ \sum_{i=1}^r\init_\sigma(I_{w_i})$ is squarefree with Stanley-Reisner complex 
$ \bigcap\{ \Delta(\init_\sigma(I_{w_i})):i \in [r]\}$
and that 
$\bigcap\{\init_\sigma(I_u):u\in \perm(A)\}$
is also squarefree with Stanley-Reisner complex 
$ \bigcup\{\Delta(\init_\sigma(I_u)):u\in \perm(A)\}$. 

By \cite[Proposition~4.8]{Wei}, \[\bigcap_{i \in [r]} \Delta(\init_\sigma(I_{w_i}))=\bigcup_{u\in \perm(A)}\Delta(\init_\sigma(I_u)).\]  Thus, all containments in Equation~\ref{eq:idealcontainments} are actually equalities.



\noindent (ii) Each ASM ideal is homogeneous with respect to the standard grading on $R$.
Because
\[\sum_{i=1}^r I_{w_i}\subseteq I_A\subseteq \bigcap_{u\in \perm(A)} I_u \mbox{ \hspace{1cm} and \hspace{1cm} } \displaystyle \init_\sigma\left(\sum_{i=1}^r I_{w_i}\right)=\init_\sigma\left (\bigcap_{u\in \perm(A)} I_u\right),
\]we know, by comparing Hilbert functions (see Subsection \href{https://arxiv.org/abs/2108.08370v4}{2.4}),
\[\sum_{i=1}^r I_{w_i}= I_A= \bigcap_{u\in \perm(A)} I_u.\]

\noindent (iii) First note that there exists a set $\{\pi_1,\ldots,\pi_m\}$ of bigrassmannian permutations in $S_n$ so that $A=\vee \{\pi_1,\ldots,\pi_m\}$ (this follows from \cite{LS96} and \cite[Proposition~9]{Rea02}). Thus, as a consequence of \ref{lempart:addIwi}, $\init_\sigma(I_A)$ is radical and so $I_A$ is also radical.

\noindent (iv) Again, noting that we can write $A= \vee \{\pi_1,\ldots,\pi_m\}$ for some set of bigrassmannian permutations, applying the argument in \ref{lempart:addIwi}, we conclude 
$ I_A= \bigcap\{ I_u:u\in \perm(A)\}$.

By Proposition~\ref{prop:matSchubFacts}, $I_u$ is prime for each $u\in \perm(A)$.  By definition, the elements of $\perm(A)$ are pairwise incomparable in Bruhat order.  Thus, applying \cite[Lemma~15.19]{MS04}, no $I_u$ properly contains any $I_{u'}$ for $u,u' \in \perm(A)$.  

\noindent (v)  We have $\hgt (I_A)=\min\{\hgt (I_u):u\in \perm(A)\}=\min\{\ell(u):u\in \perm(A)\}=\hgt(A)$.

\noindent (vi) This is immediate from \ref{lempart:asmDecomp} and \ref{lempart:asmHeight}.
\end{proof}

\begin{thebibliography}{99}
\bibitem[Ful92]{Ful92} William Fulton, \emph{Flags, {S}chubert polynomials, degeneracy loci, and determinantal formulas}, Duke Math. J. \textbf{65} (1992), no.~3, 381--420.
\bibitem[KM05]{KM05} Allen Knutson and Ezra Miller, \emph{Gr{\"o}bner geometry of {S}chubert polynomials}, Ann. of Math. (2005), 1245--1318.
\bibitem[LS96]{LS96} \bysame, \emph{Treillis et bases des groupes de {C}oxeter}, Electron. J. Combin. \textbf{3} (1996), no.~2, Research paper 27, approx. 35.
\bibitem[MS04]{MS04} Ezra Miller and Bernd Sturmfels, \emph{Combinatorial commutative algebra}, vol. 227, Springer Science \& Business Media, 2004.
\bibitem[Rea02]{Rea02} Nathan Reading, \emph{Order dimension, strong {B}ruhat order and lattice properties for posets}, Order \textbf{19} (2002), no.~1, 73--100.
\bibitem[Wei17]{Wei} Anna Weigandt, \emph{Prism tableaux for alternating sign matrix varieties}, Preprint (2017), 33 pages, {\sf arXiv:1708.07236}.
\end{thebibliography}
\end{document}
